\documentclass[]{article}

\usepackage{amsmath}
\usepackage{multirow}
\usepackage{natbib}
\usepackage{graphicx} 


\begin{document}

% Auto-generated figures file
% \usepackage{graphicx}

\begin{figure}[h!]
    \centering
    \includegraphics[width=0.5\textwidth]{../input_files/plots/step_1_wiener_filter_diagnostic_1_1775607926.png}
    \caption{Angular power spectra, $C\_{\ell}$, of the signal and noise components in the simulated map. The plot illustrates that the kinetic Sunyaev-Zel'dovich (kSZ) signal (green) is sub-dominant by several orders of magnitude compared to the primary Cosmic Microwave Background (CMB, blue), which is shown convolved with the 1.4 arcmin instrumental beam, and the instrumental white noise (red). The transfer function of the Wiener filter, $W(\ell)$ (purple, right axis), is constructed from these components to optimally suppress the dominant CMB and noise contributions by down-weighting the Fourier modes where their power is highest. This filtering is an essential step to isolate the faint kSZ signal for subsequent analysis.}
    \label{fig:image0}
\end{figure}

\begin{figure}[h!]
    \centering
    \includegraphics[width=0.5\textwidth]{../input_files/plots/step_4_ksz_analysis_summary_2_1775608915.png}
    \caption{Summary of the kinetic Sunyaev-Zel'dovich (kSZ) analysis. (a) The recovered Thomson optical depth ($\tau$) as a function of halo mass ($M\_{500}$). The measurements, shown with large Jackknife-derived error bars, are dominated by noise, yielding a weak constraint on the power-law slope that is statistically consistent with the theoretical model (dashed line). (b) A comparison of reconstructed and true line-of-sight velocities reveals no correlation (Pearson $r=0.026$), a failure attributed to the sparse halo catalog. (c) Velocity-weighted stacked cutouts of the kSZ signal for three mass bins, illustrating the low signal-to-noise of the detection. (d) The detection signal-to-noise ratio (SNR) improves from 0.41 to 1.06 after applying a Wiener filter to mitigate CMB foregrounds, but degrades when using the poorly reconstructed velocities (orange bar) instead of the true velocities (blue bar).}
    \label{fig:image1}
\end{figure}



\end{document}
                